% GENERATED FILE. Edit canonical lessons, then run npm run generate:books.
% canonical-source-hash: fnv1a64:5d783215370f907a
% canonical-lessons: TE-C363-tirugubatu, TE-C363-tiragabadu, TE-C363-anaciveyu, TE-C363-usigolpu, TE-C363-svatantryam

\chapter{Rebellion}
\label{ch:te-rebellion}
\hlchaptermodality{363}

\begin{chapteropening}
\textbf{By the end of this chapter:} \emph{I can say a rebellion, to rise up, to suppress, to incite and independence.}

Complete the last lesson of chapter 363: I can say a rebellion, to rise up, to suppress, to incite and independence.
\end{chapteropening}

\section[\texorpdfstring{\te{తిరుగుబాటు} (tirugubāṭu)}{తిరుగుబాటు (tirugubāṭu)}]{\te{తిరుగుబాటు} (tirugubāṭu) — a rebellion; a revolt}

\label{lesson:TE-C363-tirugubatu}



\begin{quote}
Before the new one: say the Telugu for a truce, then the Telugu for an envoy.
\end{quote}

\subsection*{You'll want to know: \te{తిరుగుబాటు}}
\textbf{\te{తిరుగుబాటు}} — \emph{tirugubāṭu} — \textquotedblleft{}a rebellion; a revolt\textquotedblright{}.

\begin{quote}
\te{రైతులు} \te{రాజు} \te{మీద} \te{తిరుగుబాటు} \te{చేశారు}.
\end{quote}

\begin{quote}
\emph{raitulu rāju mīda tirugubāṭu cēśāru.}
\end{quote}

\begin{quote}
The farmers rebelled against the king.
\end{quote}

\te{తిరుగుబాటు} \te{చేయు}, to rebel; from \te{తిరుగు}, to turn.

\begin{cousinweb}
\textbf{\te{తిరుగుబాటు}} (\emph{tirugubāṭu}) is native Dravidian, built on \emph{tirugu}, \textquotedblleft{}to turn, turn back\textquotedblright{} (DEDR 3246): a turning against.
\end{cousinweb}

\begin{grammarlens}[title={What you've built}]
One more word for telling what happened.
\end{grammarlens}

\subsection*{Your turn}
\begin{itemize}
\raggedright
  \item \emph{Say it:} \emph{tirugubāṭu}
  \item \emph{Say it:} \emph{tirugubāṭu}, once more
  \item \emph{Say it:} \emph{raitulu rāju mīda tirugubāṭu cēśāru.} — The farmers rebelled against the king.
  \item \emph{Recall:} say the Telugu for a truce, then the Telugu for an envoy, then say \emph{tirugubāṭu} again
\end{itemize}

\begin{tcolorbox}[breakable,colback=teal!4,colframe=teal!35!black,title={Before you move on}]
What does \te{తిరుగుబాటు} mean? (\textquotedblleft{}A rebellion; a revolt\textquotedblright{}.) Say it once more.
\end{tcolorbox}
\section[\texorpdfstring{\te{తిరగబడు} (tiragabaḍu)}{తిరగబడు (tiragabaḍu)}]{\te{తిరగబడు} (tiragabaḍu) — to rise up; to revolt}

\label{lesson:TE-C363-tiragabadu}



\begin{quote}
Before the new one: say the Telugu for an envoy, then the Telugu for a rebellion.
\end{quote}

\subsection*{You'll want to know: \te{తిరగబడు}}
\textbf{\te{తిరగబడు}} — \emph{tiragabaḍu} — \textquotedblleft{}to rise up; to revolt\textquotedblright{}.

\begin{quote}
\te{కష్టాలు} \te{భరించలేక} \te{రైతులు} \te{తిరగబడ్డారు}.
\end{quote}

\begin{quote}
\emph{kaṣṭālu bhariñcalēka raitulu tiragabaḍḍāru.}
\end{quote}

\begin{quote}
Unable to bear their hardships, the farmers rose up.
\end{quote}

Past (I …): \te{తిరగబడ్డాను}, \emph{tiragabaḍḍānu}.

Close to \te{తిరుగుబాటు} \te{చేయు}; the one opposed takes \te{మీద}.

\begin{cousinweb}
\textbf{\te{తిరగబడు}} (\emph{tiragabaḍu}) is native Dravidian: \emph{tirugu}, \textquotedblleft{}to turn\textquotedblright{} (DEDR 3246), with the helping verb \emph{paḍu}.
\end{cousinweb}

\begin{grammarlens}[title={What you've built}]
One more word for telling what happened.
\end{grammarlens}

\subsection*{Your turn}
\begin{itemize}
\raggedright
  \item \emph{Say it:} \emph{tiragabaḍu}
  \item \emph{Say it:} \emph{tiragabaḍu}, once more
  \item \emph{Say it:} \emph{kaṣṭālu bhariñcalēka raitulu tiragabaḍḍāru.} — Unable to bear their hardships, the farmers rose up.
  \item \emph{Recall:} say the Telugu for an envoy, then the Telugu for a rebellion, then say \emph{tiragabaḍu} again
\end{itemize}

\begin{tcolorbox}[breakable,colback=teal!4,colframe=teal!35!black,title={Before you move on}]
What does \te{తిరగబడు} mean? (\textquotedblleft{}To rise up; to revolt\textquotedblright{}.) Say it once more.
\end{tcolorbox}
\section[\texorpdfstring{\te{అణచివేయు} (aṇacivēyu)}{అణచివేయు (aṇacivēyu)}]{\te{అణచివేయు} (aṇacivēyu) — to suppress; to crush}

\label{lesson:TE-C363-anaciveyu}



\begin{quote}
Before the new one: say the Telugu for a rebellion, then the Telugu for to rise up.
\end{quote}

\subsection*{You'll want to know: \te{అణచివేయు}}
\textbf{\te{అణచివేయు}} — \emph{aṇacivēyu} — \textquotedblleft{}to suppress; to crush\textquotedblright{}.

\begin{quote}
\te{రాజు} \te{ఆ} \te{తిరుగుబాటును} \te{అణచివేశాడు}.
\end{quote}

\begin{quote}
\emph{rāju ā tirugubāṭunu aṇacivēśāḍu.}
\end{quote}

\begin{quote}
The king crushed that rebellion.
\end{quote}

Past (I …): \te{అణచివేశాను}, \emph{aṇacivēśānu}.

\te{అణచు}, to press down, with \te{వేయు}, which makes it final.

\begin{cousinweb}
\textbf{\te{అణచివేయు}} (\emph{aṇacivēyu}) is native Dravidian: DEDR 112 lists Telugu \emph{aṇãcu}, \textquotedblleft{}to suppress, humble, subject\textquotedblright{}; here it takes \emph{vēyu}, \textquotedblleft{}to throw\textquotedblright{}, which makes the action complete.
\end{cousinweb}

\begin{grammarlens}[title={What you've built}]
One more word for telling what happened.
\end{grammarlens}

\subsection*{Your turn}
\begin{itemize}
\raggedright
  \item \emph{Say it:} \emph{aṇacivēyu}
  \item \emph{Say it:} \emph{aṇacivēyu}, once more
  \item \emph{Say it:} \emph{rāju ā tirugubāṭunu aṇacivēśāḍu.} — The king crushed that rebellion.
  \item \emph{Recall:} say the Telugu for a rebellion, then the Telugu for to rise up, then say \emph{aṇacivēyu} again
\end{itemize}

\begin{tcolorbox}[breakable,colback=teal!4,colframe=teal!35!black,title={Before you move on}]
What does \te{అణచివేయు} mean? (\textquotedblleft{}To suppress; to crush\textquotedblright{}.) Say it once more.
\end{tcolorbox}
\section[\texorpdfstring{\te{ఉసిగొల్పు} (usigolpu)}{ఉసిగొల్పు (usigolpu)}]{\te{ఉసిగొల్పు} (usigolpu) — to incite; to egg on}

\label{lesson:TE-C363-usigolpu}



\begin{quote}
Before the new one: say the Telugu for to rise up, then the Telugu for to suppress.
\end{quote}

\subsection*{You'll want to know: \te{ఉసిగొల్పు}}
\textbf{\te{ఉసిగొల్పు}} — \emph{usigolpu} — \textquotedblleft{}to incite; to egg on\textquotedblright{}.

\begin{quote}
\te{ఎవరో} \te{ప్రజలను} \te{గొడవకు} \te{ఉసిగొల్పారు}.
\end{quote}

\begin{quote}
\emph{evarō prajalanu goḍavaku usigolpāru.}
\end{quote}

\begin{quote}
Someone incited the people to a quarrel.
\end{quote}

Past (I …): \te{ఉసిగొల్పాను}, \emph{usigolpānu}.

Whom takes -\te{ను}, what to takes -\te{కు}: \te{గొడవకు} \te{ఉసిగొల్పు}.

\begin{cousinweb}
\textbf{\te{ఉసిగొల్పు}} (\emph{usigolpu}) is a Telugu verb for setting someone on; its origin is not settled.
\end{cousinweb}

\begin{grammarlens}[title={What you've built}]
One more word for telling what happened.
\end{grammarlens}

\subsection*{Your turn}
\begin{itemize}
\raggedright
  \item \emph{Say it:} \emph{usigolpu}
  \item \emph{Say it:} \emph{usigolpu}, once more
  \item \emph{Say it:} \emph{evarō prajalanu goḍavaku usigolpāru.} — Someone incited the people to a quarrel.
  \item \emph{Recall:} say the Telugu for to rise up, then the Telugu for to suppress, then say \emph{usigolpu} again
\end{itemize}

\begin{tcolorbox}[breakable,colback=teal!4,colframe=teal!35!black,title={Before you move on}]
What does \te{ఉసిగొల్పు} mean? (\textquotedblleft{}To incite; to egg on\textquotedblright{}.) Say it once more.
\end{tcolorbox}
\section[\texorpdfstring{\te{స్వాతంత్ర్యం} (svātantryaṁ)}{స్వాతంత్ర్యం (svātantryaṁ)}]{\te{స్వాతంత్ర్యం} (svātantryaṁ) — independence; freedom}

\label{lesson:TE-C363-svatantryam}



\begin{quote}
Before the new one: say the Telugu for to suppress, then the Telugu for to incite.
\end{quote}

\subsection*{You'll want to know: \te{స్వాతంత్ర్యం}}
\textbf{\te{స్వాతంత్ర్యం}} — \emph{svātantryaṁ} — \textquotedblleft{}independence; freedom\textquotedblright{}.

\begin{quote}
\te{మన} \te{దేశానికి} \te{స్వాతంత్ర్యం} \te{ఎప్పుడు} \te{వచ్చింది}?
\end{quote}

\begin{quote}
\emph{mana dēśāniki svātantryaṁ eppuḍu vaccindi?}
\end{quote}

\begin{quote}
When did our country gain independence?
\end{quote}

\te{స్వ}, self, and \te{తంత్రం}, rule: self-rule.

\begin{cousinweb}
\textbf{\te{స్వాతంత్ర్యం}} (\emph{svātantryaṁ}) is borrowed from Sanskrit \emph{svātantrya}, \textquotedblleft{}following one's own will, independence\textquotedblright{} (Monier-Williams), from \emph{sva-tantra}, \textquotedblleft{}self-ruled\textquotedblright{}.
\end{cousinweb}

\begin{grammarlens}[title={What you've built}]
One more word for telling what happened.
\end{grammarlens}

\subsection*{Your turn}
\begin{itemize}
\raggedright
  \item \emph{Say it:} \emph{svātantryaṁ}
  \item \emph{Say it:} \emph{svātantryaṁ}, once more
  \item \emph{Say it:} \emph{mana dēśāniki svātantryaṁ eppuḍu vaccindi?} — When did our country gain independence?
  \item \emph{Recall:} say the Telugu for to suppress, then the Telugu for to incite, then say \emph{svātantryaṁ} again
\end{itemize}

\begin{tcolorbox}[breakable,colback=teal!4,colframe=teal!35!black,title={Before you move on}]
What does \te{స్వాతంత్ర్యం} mean? (\textquotedblleft{}Independence; freedom\textquotedblright{}.) Say it once more.
\end{tcolorbox}
